% =============================================================================
% Section 08 — Invariance and Continuity
% Book I — Ontology of Viability
% Chapter 02 — Invariance
%
% Purpose:
%   To establish the ontological distinctions between invariance,
%   continuity, and identity, and to examine their possible
%   relations under transformation.
%
% Conceptual Policy:
%   1. Invariance is not identical to continuity.
%   2. Continuity is not identical to numerical identity.
%   3. The absence of identified invariants does not establish
%      discontinuity or the loss of identity.
%   4. The preservation of invariants does not independently
%      establish continuity or numerical identity.
%   5. No universal criterion of continuity is assumed.
% =============================================================================

\section{Invariance and Continuity}
\label{sec:invariance-and-continuity}

\subsection{The Ontological Problem of Continuity}

Continuity constitutes a distinct ontological problem within
the analysis of transformation. It concerns the manner in which
an existent, or a determinate process involving an existent,
is connected across successive states, phases, or transformations.

The concept of continuity is frequently associated with the
preservation of certain determinations. When an existent changes
while some of its characteristics remain preserved, this
preservation may be interpreted as evidence of continuity.

However, such an interpretation requires further examination.

The preservation of a determination and the continuity of
an existent are not identical questions. A determination may
remain preserved across a transformation without establishing
that the existent itself has continued. Conversely, an existent
may undergo extensive transformation without the absence
of a particular preserved determination establishing a
discontinuity in its being.

Continuity must therefore be distinguished from invariance.

Invariance concerns the preservation of a specified
determination under a specified transformation. Continuity
concerns a relation of connectedness across successive
states or phases.

Although these concepts may be related, neither can be
reduced to the other.

\subsection{Continuity as a Relation Across Transformation}

Continuity is not simply the absence of change.

An existent may undergo continuous transformation while
its determinations change. The occurrence of transformation
does not necessarily interrupt continuity, just as the
absence of an immediately identifiable transformation does
not independently establish it.

The analysis of continuity therefore requires attention
to the relation between successive states or phases.

Such a relation may involve:

\begin{itemize}
    \item an ordered succession of states;
    \item a connection between successive phases;
    \item the persistence of a process through transformation;
    \item a relation of dependence between earlier and later
          states;
    \item a determinate organization of transformation.
\end{itemize}

These aspects represent possible dimensions of continuity.
They should not be treated as universally necessary conditions
for every form of continuity.

Different existents and processes may exhibit different
modes of connectedness. Consequently, continuity cannot
be defined through one undifferentiated criterion applicable
to every ontological domain.

The relevant relation must be specified according to
the existent or process under consideration.

\subsection{The Distinction Between Invariance and Continuity}

Invariance and continuity answer different analytical questions.

Invariance asks whether a specified determination remains
preserved under a specified transformation.

Continuity asks whether successive states or phases stand
in a determinate relation of connectedness.

The distinction becomes clear when their possible combinations
are considered.

A determination may remain invariant across successive
states, while the relation between the existents occupying
those states remains undetermined.

Likewise, an existent may undergo continuous transformation
even when none of its previously identified determinations
remains invariant.

Thus, invariance does not automatically establish continuity,
and continuity does not necessarily require the preservation
of a particular invariant.

The two concepts may be associated in particular cases,
but their relationship must be established rather than assumed.

\subsection{Continuity Without Identified Invariants}

The possibility of continuity without an identified invariant
is philosophically significant.

An existent may undergo a transformation in which its
previously identified determinations are altered. The absence
of a preserved determination does not, by itself, establish
that the existent has become discontinuous.

This conclusion follows from the distinction between
the preservation of determinations and the connectedness
of successive states.

A transformation may involve extensive alteration while
maintaining a determinate relation between successive
phases. The relevant continuity may therefore concern
the organization or connection of the transformation
rather than the preservation of a particular feature.

It is also possible that an invariant exists but has not
been identified. However, the ontological possibility
of continuity must not be made dependent upon the successful
identification of such an invariant.

The absence of an identified invariant is an epistemically
limited statement about the available determination of
preservation. It is not, by itself, an ontological proof
of discontinuity.

\subsection{Invariance Without Established Continuity}

The converse possibility must also be considered.

A specified determination may remain preserved across
different states or transformations without establishing
that the existents involved constitute one continuous
existent.

For example, distinct existents may possess the same
structural, compositional, or quantitative determination.
The preservation of that determination across them does
not establish a continuous passage from one existent
to another.

Similarly, an invariant may be preserved across a sequence
of states whose relation to one another has not been
ontologically established.

In such cases, invariance concerns the repeated preservation
of a determination, whereas continuity remains a separate
question concerning the connectedness of the states or
phases.

Consequently, the presence of invariants cannot serve
as an independent proof of continuity.

\subsection{Continuity and Numerical Identity}

Continuity must also be distinguished from numerical identity.

Numerical identity concerns whether an existent is one
and the same existent. Continuity concerns a relation
of connectedness across states or phases.

These concepts may be related, particularly when the
persistence of an existent through transformation is
under consideration. Nevertheless, continuity cannot
be treated as a sufficient definition of numerical identity.

A continuous process may involve distinct existents
at different stages. Likewise, a relation of continuity
between successive states does not, without further
ontological determination, establish that those states
belong to one numerically identical existent.

Conversely, the analysis of numerical identity cannot
be replaced by an analysis of continuity alone.

To infer numerical identity from continuity would require
an additional account of why the connected states belong
to one and the same existent rather than merely participating
in a continuous process or relation.

The distinction therefore remains necessary even where
continuity and identity appear closely associated.

\subsection{Continuity and Persistence}

Continuity is also distinct from persistence.

Persistence concerns the continued existence of an existent
across a temporal or transformational domain. Continuity
concerns the connectedness of successive states or phases.

Although persistence may involve continuity, the two
concepts are not interchangeable.

An existent may persist through transformations whose
successive phases are not adequately described by a simple
model of continuity. Conversely, a continuous process
may extend across a succession of distinct existents
without establishing the persistence of one numerically
identical existent.

The distinction prevents the assumption that every
continuous sequence constitutes the persistence of
one existent, or that every form of persistence must
be described through the same model of continuity.

The conditions under which persistence and continuity
coincide must therefore be examined separately.

\subsection{Continuity and Stability}

Continuity must not be confused with stability.

Stability concerns the maintenance of a condition,
configuration, or mode of organization under specified
circumstances or perturbations.

Continuity concerns connectedness across successive
states or phases.

A process may remain continuous while undergoing unstable
or rapidly changing configurations. Conversely, a stable
condition may be observed at distinct moments without
the stability itself establishing a continuous relation
between them.

Stability may contribute to the preservation of particular
determinations, but it is neither identical to continuity
nor a universal requirement for it.

The distinction is important because continuity should
not be interpreted as the persistence of an unchanged
condition.

\subsection{The Transformation of Continuity}

Continuity itself may be subject to different forms of
transformation.

The relation between successive states may remain preserved
while the determinations of those states change. It may
also undergo modification, interruption, or reorganization.

These possibilities require distinguishing transformation
within continuity from transformation of the relation
that constitutes continuity.

A change occurring within a continuous process does not
necessarily alter the continuity of that process. However,
a transformation affecting the relation between successive
phases may raise a distinct question concerning whether
the previous continuity remains applicable.

No universal conclusion can be derived merely from the
magnitude or number of changes involved.

The ontological significance of a transformation depends
on the particular relation of connectedness under
consideration.

\subsection{The Limits of Invariance as a Criterion of Continuity}

The preceding distinctions establish the limits of
using invariance as a criterion of continuity.

The preservation of an invariant may provide a relevant
determination when examining a continuous transformation.
Nevertheless, it cannot independently establish that
continuity exists.

Likewise, the failure to identify an invariant cannot
independently establish that continuity has ceased.

Any argument concerning continuity must therefore
identify the relation connecting the relevant states
or phases rather than relying exclusively on the
preservation or alteration of particular determinations.

This requirement prevents invariance from being assigned
an explanatory role beyond its proper analytical scope.

\subsection{Synthesis}

The analysis establishes the following distinctions:

\begin{itemize}
    \item Invariance concerns the preservation of a specified
          determination across transformation.
    \item Continuity concerns the connectedness of successive
          states or phases.
    \item Persistence concerns the continued existence of
          an existent across a specified domain.
    \item Numerical identity concerns whether an existent
          is one and the same existent.
    \item Stability concerns the maintenance of a condition
          or organization under specified circumstances.
\end{itemize}

These concepts may be related, but none can be substituted
for another without an additional ontological argument.

Continuity may occur without an identified invariant,
and invariance may be observed without established continuity.
Neither possibility determines numerical identity by itself.

The preservation of invariants is therefore one possible
aspect of the analysis of continuity, not its universal
foundation or sufficient criterion.

Continuity remains a distinct ontological problem whose
relation to identity, persistence, and transformation
must be examined without reducing it to the preservation
of unchanged determinations.

% =============================================================================
% End of Section 08
% =============================================================================